\section{Secondo Piatto}

\begin{quotex}
I never said or thought that different traditions had to or even could merge or assimilate in any way whatsoever, even in their theoretical part. All that is possible and desirable, is a common understanding on a certain ground. This cannot be disputed when one admits that, the truth being one, it is possible to establish equivalences between its various modes of expression. I would further add that communication between various traditional doctrines, under conditions which do not detract from their independence, can only be effectuated from above. Perhaps you have already said that; unfortunately it is not possible for me currently to explain myself as clearly as I would like, especially as it would be very difficult, I fear, to find qualified individuals capable of taking the initiative of an effective rapprochement such as the one I am thinking of. Ultimately, this may come. Never despair.
\flright{\textit{Letter from Rene Guenon to Noëlle Maurice-Denis}, 30 March 1919}
\end{quotex}


For the second phase of Gornahoor, I have listened to comments and suggestions. So instead of reprising the same topic in different contexts, I will have a few core topics which will be updated and corrected as necessary. They will be the foundation for everything else.

The twin purposes of Gornahoor are these:

\begin{itemize}


\item Restate traditional metaphysics in ``western'' terms. I use the word in a broad sense to include all areas that were, at one time or another, part of the Roman Empire, which is not at all identical to ``European''.

\item Update the Path of the Knight for a contemporary context. This is different from exoteric paths or the monk's way. It is for men of action, engaged in the world in the search for the Holy Grail. This is the same as the Sophianic path.
\end{itemize}


\paragraph{Exoteric and Esoteric Spirituality}


\begin{quotex}
Jews ask for signs, and Greeks look for wisdom.
\flright{\textsc{1 Colossians 1:22}}
\end{quotex}


Consciousness has a history just like any other human manifestation. Esoterically the above quote does not represent two ethnic groups, but serve as types for the historical manifestation of spiritual understanding.

\begin{itemize}


\item \textbf{Exoteric}: At this level, religious phenomena are experienced as wondrous events, experienceable by the senses.

\item \textbf{Rational}: Here, extraordinary events are rejected and the physical sciences serve as the model to understand the world. The cost of this, unfortunately, is quite high and requires the rejection not only of free will but of consciousness itself.

\item \textbf{Esoteric}: At the stage, religious experience is understood by the higher intellect, which regards the senses and mere rational thinking as inadequate.
\end{itemize}


Keep in mind that the lower cannot understand the higher.

\paragraph{Unity of Metaphysics}


I don't like the phrase ``Transcendent Unity of Religions'' since it is unclear and even misleading. Instead, our goal is more precise: To show the compatibility of the metaphysics assumed by Traditional religions. Hence I shall rely, for the most part, on these three sources:

\begin{enumerate}


\item \emph{Man and his Becoming} by \textbf{Rene Guenon}


\item \emph{An Introduction to Philosophy} by \textbf{Jacques Maritain}


\item \emph{Prolegomena to the Metaphysics of Islam} by \textbf{Syed Muhammad Naquid al-Attas}
\end{enumerate}


I have no intention to repeat their arguments, but will instead accept their conclusions as established. This task will require some compromises from the three sources.

\paragraph{Forerunners}


At the first half of the twentieth century, there are several thinkers who have changed our understanding of religious experience. Obviously, this is a personal list, which I have categorized in this way:

\begin{itemize}


\item \textbf{Sophia}: Vladimir Soloviev, Sergei Bulgakov, Nikolai Berdyaev, Henry Corbin

\item \textbf{Tradition}: Rene Guenon, Ananda Coomaraswamy, Frithjof Schuon, Mircea Eliade, D T Suzuki

\item \textbf{Psychology}: Carl Jung (qualified), James Hillman (qualified)

\item \textbf{Hermetism}: Valentin Tomberg, Boris Mouravieff

\end{itemize}


I will justify these choices along the way. There are three realms: the material, the imaginal, and the celestial. If the emphasis is solely on the celestial, it can be difficult to relate it to ordinary experience.

\paragraph{Degree of Difficulty}


This is a difficult path; after all, only 3 Knights found the Holy Grail out of the 150 who started. This is meant for those three. Failure comes from this misunderstanding:

\textbf{Tradition is a way of life, not a system of thought.}


\subsection*{Upcoming Topics}


\begin{quotex}
Everything must be conscious, at least potentially, or it does not exist

Anything said or any event must have meaning, otherwise it may as well never have happened.
\end{quotex}


This is a tentative list of the core topics. Other posts will be inserted from time to time, which may include translations, comments on current events, book and movie reviews, and even recipes. With some exceptions, I will use the word ``philosophy'' as a more general term for traditional metaphysics, not as profane philosophy in the academic sense.

\paragraph{Bathsheba and the Queen of Heaven}


The story of Bathsheba will serve as the foundation for the path of Sophia.

\paragraph{Eternal Sophia}


A review of Henry Corbin's essay on Jung's \emph{Answer to Job}, as a revelation of Sophia.

\paragraph{First Principles}


First principles, common sense, and the divisions of Philosophy

\paragraph{Metaphysical Trinity}


Vladimir Soloviev, in \emph{Lectures on Divine Humanity}, casually mentions that the idea of the trinity arises from Plotinus. Although it is metaphysically true, its religious and theological expressions are not so simple.

\paragraph{The World Within}


The world has an inside as well as an outside.

\paragraph{Three Gunas}


The three forces in the world: positive, negative, and reconciling. The last one can be the most difficult one to discern.

\paragraph{Intellectual Love of God}


The highest love of God is in the intellect, not in the senses or emotions.

\paragraph{Purpose of Life}


The purpose of life is not created.

\paragraph{Contemporary Philosophy}


These are personal interests from academic philosophy, which are not devoid of interest: Nagel's bat, the Chinese box, the Turing Test, and possible worlds.

\paragraph{Living in a Simulation}


How to determine if we are living in a simulation or a real world.

\paragraph{The Three Body Problem}


This problem has no computable solution, yet everything moves as expected.

\paragraph{Being and Nonbeing}


Some light on nonbeing.

\paragraph{Being and Existence}


Being: that which is; Existence; the act of being.

\paragraph{Substance and Accident}


What is essential to a being and what is not.

\paragraph{Potentiality and Act}


Actualization of Being

\paragraph{The Person}


What is a person?

\paragraph{The Possible and the Real}


Their metaphysical equivalence.

\paragraph{Degrees of Existence}


Different states of existence and their relationships.

\paragraph{Descent of Man}


How the human being descends through the states of being

\paragraph{History of Consciousness}


From naïve participation, to rationalism, and so on.

\paragraph{Great Being}


Auguste Comte meets Vladimir Soloviev in a strange but intriguing way.

\paragraph{Saving the Appearances}


The nature of scientific knowledge.

\paragraph{Unity of Being}


If Being is the principle of beings, how can there be individual beings?

\paragraph{Justice and Mercy}


Justice and the works of mercy.

\paragraph{The Common Task and Cosmism: Reflections on Nikolai Fyodorovich Fyodorov.}


\flrightit{Posted on 2021-08-22 by Cologero}
